\section{Related Work}

\subsection{Automated Leukocyte Morphology Classification}

Automated blood-cell image analysis has a long history, but recent progress has been driven by annotated single-cell datasets and convolutional neural networks. Reviews of white blood cell classification report a broad literature on segmentation, feature extraction, transfer learning, and deep classifiers for microscopic blood-smear imagery \cite{khan2021_wbc_review,deshpande2021_blood_cells_review}. Much of this work evaluates a fixed set of classes and reports accuracy, precision, recall, or F1 score under same-dataset splits. Such studies are useful for establishing classification feasibility, but they generally do not test whether a model can refuse a sample outside its label inventory.

Several influential hematology image-analysis studies demonstrate the strength of deep learning under closed-set or task-specific settings. Acevedo et al. introduced a peripheral blood cell dataset and reported convolutional recognition of multiple cell categories \cite{acevedo2019_peripheral_blood_cnn}. Matek et al. showed high-performance blast-cell recognition in acute myeloid leukemia images and later extended deep learning to large bone marrow morphology classification \cite{matek2019_blast,matek2021_bone_marrow}. Walter et al. review the broader promise and limitations of artificial intelligence in hematological diagnostics \cite{walter2023_ai_hematology}. These studies establish the relevance of automated morphology analysis but do not remove the need for open-set evaluation. Strong classification of known cell types is a different claim from reliable rejection of unseen morphologies.

This distinction matters because morphology labels are not interchangeable with patient-level clinical endpoints. A smear image can be assigned a cell-type or maturation label without determining the patient's diagnosis, and a model can be useful for one of these tasks while being inappropriate for the other. The present manuscript therefore keeps the task at the single-cell morphology level. Its central question is not whether a patient has leukemia or whether a sample is clinically abnormal, but whether a classifier trained on mature leukocytes can recognize when an image belongs outside that fixed label inventory. That framing follows the data available in the frozen artifacts and avoids translating morphology rejection into a diagnostic claim that the experiments do not test.

The MLL23 dataset is particularly relevant because it contains more than 40,000 expert-annotated single-cell images in 18 morphologically defined classes, with 288 by 288 pixel TIFF images after curation \cite{mll23_nature_2025,mll23_zenodo_2024}. Its class diversity makes it possible to define an explicit open-set protocol by training on mature leukocytes and holding out other named morphologies. The AML-Cytomorphology\_LMU collection is also central because it provides an independent TCIA source with 18,365 expert-labeled single-cell images from peripheral blood smears \cite{aml_lmu_tcia_2019}. The combination of MLL23 and AML-LMU allows the paper to ask not only whether unknown morphologies can be rejected internally, but whether the same conclusions survive external acquisition differences.

\subsection{Open-Set Recognition and Out-of-Distribution Detection}

Open-set recognition was formalized to address settings where not all test classes are known during training \cite{scheirer2013_open_set}. In deep learning, OpenMax showed that closed-set neural networks can assign high confidence to inputs outside the training taxonomy and introduced a recalibrated output layer for unknown rejection \cite{bendale2016_openmax}. Surveys of open-set recognition emphasize that the task requires balancing known-class recognition with unknown-class rejection, and that evaluation protocols must avoid using unknown-test data during method selection \cite{geng2021_osr_survey}.

Out-of-distribution detection is closely related but not identical. Generalized OOD surveys unify anomaly detection, novelty detection, open-set recognition, OOD detection, and outlier detection while highlighting differences in training assumptions and semantic shift \cite{yang2024_generalized_ood}. In medical image analysis, OOD detection is especially important because models may encounter unseen pathologies, acquisition changes, artifacts, or population differences \cite{hong2024_medical_ood_survey}. The present study lies at the overlap of open-set recognition and medical OOD detection: the unknown samples are semantic morphology classes held out from training, while external AML-LMU evaluation also introduces acquisition and taxonomy-domain differences.

Several post-hoc OOD scores are relevant baselines. Maximum softmax probability remains a strong and simple method because closed-set confidence is often informative even when imperfect \cite{hendrycks2017_msp}. Energy scores use the log-sum-exp of logits to reduce reliance on the top softmax probability alone \cite{liu2020_energy}. ReAct clips internal activations to reduce overconfident behavior on OOD inputs \cite{sun2021_react}. ViM combines a feature-space residual term with logit evidence through virtual-logit matching \cite{wang2022_vim}. These methods differ in what they fit after training, so leakage control matters. Fitting PCA subspaces, centroids, covariance estimates, or clipping thresholds with unknown-test data would invalidate an open-set evaluation.

The leakage issue is especially subtle in biomedical studies because the same external dataset can easily become both a validation resource and a reported test resource. For example, choosing a threshold, selecting a score family, deciding whether to normalize stains, or fitting a feature residual subspace after looking at external labels can improve apparent performance while changing the scientific question. Such choices may be appropriate in a deployment-engineering study, but they no longer measure zero-shot transfer of the frozen internal protocol. The present manuscript therefore treats post-hoc scoring as part of the internal pipeline and keeps AML-LMU completely outside all fitting and selection steps.

\subsection{Representation Learning and Unknown Rejection}

Representation learning can improve known-class embedding geometry. Supervised contrastive learning encourages samples from the same class to cluster while separating samples from different classes \cite{khosla2020_supcon}. ArcFace imposes an angular margin between normalized features and normalized class weights and has been highly influential in face recognition \cite{deng2019_arcface}. In principle, tighter known-class clusters and larger inter-class margins could improve open-set rejection because unknown samples might lie farther from all known classes.

That hypothesis is plausible but not guaranteed. Unknown morphologies can be visually close to known classes, and a representation that improves known-class compactness may also make near-taxonomy unknowns more confidently attracted to a known prototype. For this reason, representation-learning results must be judged by multiseed and split sensitivity, not by a single favorable run. The present manuscript treats ArcFace as a major ablation and asks whether its internal gains reproduce across seeds, splits, and external validation. It does not present ArcFace as a newly proposed hematology method.

\subsection{Domain Shift and External Validation in Medical Imaging}

Domain shift is a persistent obstacle in medical image analysis. Differences in scanners, staining, acquisition protocols, preprocessing, patient populations, and annotation conventions can all degrade model performance outside the source dataset \cite{litjens2017_medical_ai_survey,guan2022_domain_adaptation_medical}. In blood-cell morphology, cross-domain classification studies have reported that performance can drop when testing across datasets or imaging sources \cite{salehi2022_cross_domain,sadafi2024_continual_wbc,tsutsui2023_wbc_domain_shift}. Such evidence is important because many closed-set benchmarks overestimate generalization when train and test images come from the same source.

External validation is even more stringent in open-set recognition. A model may face both covariate shift among known classes and semantic shift in the unknown set. A post-hoc score that ranks known and unknown samples well internally may fail when known samples from the new domain acquire lower confidence or when unknown samples resemble known classes in a different way. This motivates the test-only AML-LMU design: no external data are used to adapt the model, tune thresholds, choose scorers, fit PCA residuals, or adjust preprocessing. The external result therefore measures the generalization of the frozen internal protocol rather than the success of external adaptation.

This choice also shapes how negative results should be interpreted. A poor external AUROC is not evidence that adaptation would be impossible, and a high FPR@95TPR is not a final clinical operating point. Instead, these results quantify how much reliability remains when a closed-set morphology classifier and its post-hoc open-set scores are transferred without external assistance. In that sense, the external experiment is closer to an audit of robustness than to an optimization benchmark. It asks whether internal model ranking is trustworthy when the data source changes, which is a practical concern for biomedical image-analysis papers that report only same-source splits.

\subsection{Topological Analysis of Learned Representations}

Persistent homology summarizes connected components, loops, and higher-dimensional features across filtrations \cite{edelsbrunner2002_persistence,carlsson2009_topology}. Computational libraries such as GUDHI support cubical and simplicial constructions for images and point clouds \cite{maria2014_gudhi}. In biomedical image analysis, topology can be used either as a feature representation or as a descriptive lens on learned embeddings. These two roles should not be conflated.

This manuscript keeps topology secondary. Cubical persistent homology is evaluated as a predictive ablation over image-derived filtrations and fusion designs. Vietoris-Rips persistent homology is used only after the predictive experiments are frozen, to describe embedding-cloud geometry across known and hard unknown morphologies. The resulting evidence is useful because it prevents overclaiming: the topology results support negative and explanatory conclusions, not a new topology-based open-set method.
